\uuid{GD7Z}
\titre{Logique}
\auteur{bodin ; croizat ; sacre}
\organisation{exo7}
\licence{CC BY-NC-SA 4.0}
\datecreate{2022-03-08}
\niveau{L1}
\module{Algèbre}
\chapitre{Logique, ensemble, raisonnement}
\sousChapitre{Logique}
\difficulte{2}
\theme{Logique}
\exo7qcmid{hilisit.logique.007}

\contenu{
\begin{question}
\qclassification{100.01}
Soit $n \ge 3$ un entier.
On considère l'implication :
$$\text{"}n \text{ nombre premier } \implies  n \text{ est impair".}$$
Quelles sont les affirmations vraies ?
\begin{answers}
	\bad{L'implication réciproque est "$n$ est pair $\implies$ $n$ est un nombre premier".}
	\good{La contraposée est "$n$ est pair $\implies$  $n$ n'est pas nombre premier".}
	\bad{Si l'implication est vraie alors l'implication réciproque l'est aussi.}
	\good{Si l'implication est vraie alors sa contraposée l'est aussi.}
\end{answers}
\begin{explanations}
L'implication réciproque est "$n$ est impair $\implies$ $n$ est un nombre premier" (ce qui est une affirmation fausse). \\
La contraposée est "$n$ est pair $\implies$  $n$ n'est pas nombre premier" (ce qui est une affirmation vraie). \\
Une implication directe peut être vraie sans que l'implication réciproque soit vraie ; c'est le cas ici.
Une implication et sa contraposée sont des propositions équivalentes (donc vraies en même temps, et fausses en même temps). \\
\end{explanations}
\end{question}
}
